%% 无标题（4 月 27 日）
\section[无标题（4 月 27 日）]{无标题的一讲（4 月 27 日）}

\srcpage{69}

今天我们要讨论非阿贝尔情形下的线算子（line operator）\yizhu{线算子指支撑在一维子流形上的算子；支撑在一点上的则叫局部算子。}{I（本库 SY02-37 术语表）}。回忆一下：若 $Z$ 是一个 4 维 TFT，则 $Z(\Sigma^2)$ 是线算子构成的范畴。它作用在配给曲面的范畴上，作用方式是沿柱状协边挖掉一个小球。算子乘积展开（operator product expansion）在 $Z(\Sigma^2)$ 上定义了一个对称幺半结构。

设 $R$ 是对称幺半高阶范畴 $\cC$ 中的一个交换代数。假定 $\cC$ 有余极限（其实我们只需要某几类），并设 $\Sigma$ 是一个拓扑空间，以组合方式给出（即一个单纯集）。于是 $\cC$ 中有一个交换代数，记作 $R\otimes\Sigma$。它由 $R^{\otimes v(\Sigma)}$（顶点）搭成：1-单形添关系，2-单形添关系之间的关系，依此类推。这个构造另有一个名字，即 $R$ 在 $\Sigma$ 上的因子化同调（factorization homology），记作 $\int_\Sigma R$；不过我们不需要这么一般的视角。\yizhu{原文把这一运算排成 $R\otimes\Sigma$，与同段随即给出的记号 $\int_\Sigma R$ 不一致，疑为排版时上标丢失。译文照排，未作改动。}{A（原文 p.69）}

在良好的情形下，$R\otimes\Sigma$ 有很好的描述。例如，若 $R$ 是交换环，则
\[
  R\otimes\Sigma \cong \cO(\mathrm{Map}(\Sigma,\Spec R)),
\]
而若 $R$ 是对称幺半范畴，则
\[
  R\otimes\Sigma \cong \QCoh(\mathrm{Map}(\Sigma,\Spec R)).
\]
这里从 $\Sigma$ 到 $\Spec R$ 的映射指局部常值映射，仍用先前的 $\pi_{\le 1}\Sigma$。所以我们为 $\Sigma$ 的每个点配一份 $R$，道路给出这些副本之间的等同。接着 2-单形给出等同之间的等同，依此类推；不过对离散环，我们看不见这些。

\srcpage{70}

上次我们还讨论了拓扑 Maxwell 理论 $\cA_{GL_1}$。它没有非平凡的局部算子，因为 $\cA_{GL_1}(S^3)$ 很无聊。然而它有非平凡的线算子，等价地说，非平凡的余 3 维缺陷。这是因为 $S^2$ 上有非平凡的主 $GL_1$-丛，而 $S^3$ 上没有。更具体地说，$\cA_{GL_1}(S^2)$ 是 $\Pic(S^2)=\Z$ 上层的范畴，于是我们得到 $\Z$-分次向量空间的范畴。这个分次是主 $U_1$-丛的 Chern 类，或该算子的单极子数。有一个到 $\C[\Z]$ 的分次维数算子，它告诉我们能在流形的结上放什么（TODO：哪个维数？）。

\begin{remark}
这推广到 3 维和 5 维：Yang–Mills 理论有有趣的奇点。这是因为 Yang–Mills 方程在 $S^2$ 上有有趣的解。于是只要指定主丛在环绕 $S^2$ 上的拓扑类型，我们就得到缺陷。这些一律称为 't Hooft 算子。
\end{remark}

而如今这些作用汇集起来：对曲面 $\Sigma$ 与 $x\in\Sigma$，$\Z$-分次向量空间的对称幺半范畴作用在范畴 $\cA_{GL_1}(\Sigma)$ 上。它们拼成一个张量积 $\Sigma\Vect_{\Z\text{-}\mathrm{gr}}$。为辨认它，我们用 Cartier 对偶：$\Z$-分次向量空间等价于 $\G_m$-表示，也就是 $\mathrm{pt}/\G_m$ 上的拟凝聚层。因此我们想知道对称幺半范畴
\[
  \Sigma\otimes\QCoh(\mathrm{pt}/\G_m) \simeq \QCoh(\mathrm{Map}(\Sigma,\mathrm{pt}/\G_m))
  \simeq \QCoh(\mathrm{Loc}_{\G_m}(\Sigma)).
\]
这有点奇怪：B 侧作用在 A 侧。我们本来预期的是另一回事，即二者等价——它们确实等价。事实上，这个作用把 $\cA_{GL_1}$ 实现为 $\mathcal{B}\G_m$ 上的自由秩 1 模。这让人想起 Fourier 变换里已经能看到的东西：带逐点乘法的 $\mathrm{Fun}(G)$ 与带卷积的 $\mathrm{Fun}(\widehat G)$ 等同，而后者通过特征标的卷积作用在前者上。

事实上，$\cB_{GL_1}(\Sigma)=\QCoh(\mathrm{Loc}_1(\Sigma))=\Sigma\otimes\cA_{GL_1}(S^2)$。由 $\cA_{GL_1}(S^2)$ 对 $\cA_{GL_1}(\Sigma)$ 的作用，诱导出 $\cB_{GL_1}(\Sigma)$ 对 $\cA_{GL_1}(\Sigma)$ 的作用；这个作用把 $\cA_{GL_1}(\Sigma)$ 实现为 $\mathcal{B}\G_m$ 上的自由秩 1 模范畴。

若 $\Sigma$ 是黎曼面，则 $\cA_{GL_1}(\Sigma)=\mathrm{Loc}(\Pic(\Sigma))$（暗地里它是 $\Pic(\Sigma)$ 上的 D-模）。我们想说单位对象是摩天层，但实际上必须取万有覆盖 $\wt{\Pic}(\Sigma)\to\Pic(\Sigma)$，它是可缩的。这大概有点奇怪。

\begin{example}
$\mathrm{Loc}(S^1)=\Rep_\Z=\QCoh(\G_m)$。$\QCoh(\G_m)$ 中的单位对象是 $\cO_1$，即单位元处的摩天层。在这个对应下，它被送到（TODO：这点我不确定）万有覆盖 $\R\to S^1$，视为一个 $\Z$-局部系统。
\end{example}

在非阿贝尔的故事里会发生类似的事；这个对象叫做 Whittaker 层。

\srcpage{71}

设 $G$ 是复约化群，例如 $\mathrm{GL}_n(\C)$、$\mathrm{SO}_n(\C)$、$\mathrm{Sp}_n(\C)$，或某个像 $E_8$ 这样的群的复形式。理论 $\cA_G$ 线性化 $\Bun_G$，并研究它的拓扑。数域类比把某个上同调自守形式的空间配给数域 $\Spec\cO_G$。类似地，对有限域 $\F_q$ 上的曲线 $C$，$\cA_G$ 理论线性化 $G$-丛的空间：$C_c[\Bun_G(C)(\F_q)]$。

$\cA_G$ 理论在物理意义下是良定义的。意思是存在一个具体构造：规范群取 $G$ 的紧形式，4 维 $\mathcal{N}=4$ 超 Yang–Mills，取 A-扭（A-twist）。见 Kapustin–Witten 的开创性论文（TODOcite）。他们没有确定闭 3-流形上的状态空间，我们也不讨论那个；我们转而聚焦于曲面。

设 $\Sigma$ 是黎曼面。我们将研究模叠 $\Bun_G(\Sigma)$，并给它配上一个 $\Bun_G(\Sigma)$ 上「拓扑」层的范畴。做法有几种，但都内蕴于拓扑，而非代数几何。

这里有一个重要之点，或许能说明我们为何对此感兴趣：我们要研究这个 TFT 中的线算子。\footnote{与 $G=GL_1$ 一样，将没有非平凡的局部算子，除非在导出意义下。在导出意义下，它们对应于一个 $G$-等变上同调环，最终对我们没有意思。}其中有许多有趣的线算子，包括 't Hooft 单极子。在 Yang–Mills 理论中，它们是余 3 维奇点。

我们首先想知道 $\cA_G(S^2)$ 是什么，这意味着要研究 $\Bun_G(\P^1)$。这由 Grothendieck–Birkhoff 定理已知。例如对 $G=\mathrm{GL}_n$，主 $\mathrm{GL}_n$-丛等价于 $\P^1$ 上秩 $n$ 的向量丛。Grothendieck–Birkhoff 定理表明 $V$ 是线丛的直和，且这个分解在排列下唯一。因此 $\Bun_n(\P^1)\simeq\Z^n/S_n$：$\mathrm{GL}_n$ 的余权格模去置换。

于是主丛很多，线算子也就很多。它们告诉你环绕 $S^2$ 内部的奇点：这些线丛的非平凡特征数，阻碍场延拓进环绕 $S^2$ 所围的球。

另一种画法是塌缩时间：做成一个意大利饺（ravioli，还是该叫 raviolo？）$\Sigma\amalg_{\Sigma\setminus x}\Sigma$，其中 $x$ 是环绕 $S^2$ 在 $\Sigma$ 内的中心。我们知道这个饺子上 $G$-丛是什么：
\[
  \Bun_G(\Sigma\amalg_{\Sigma\setminus x}\Sigma)
  \simeq \Bun_G(\Sigma)\times_{\Bun_G(\Sigma\setminus x)}\Bun_G(\Sigma).
\]
也就是说，意大利饺上的主 $G$-丛，是 $\Sigma$ 上两个丛的数据，加上 $\Sigma\setminus x$ 上的一个等同。在物理中，你可以把它想成从一个 $G$-丛 $P_1$ 出发、到另一个丛 $P_2$ 结束；我们在 $\Sigma\times[0,1]$ 的中部 $(x,1/2)$ 插入了线算子。但我们把时间塌缩成了意大利饺。因此 $P_0$ 与 $P_1$ 必须在 $x$ 之外处处一致。你可以什么都不做（$P_0=P_1$，等同取恒等），也可以做一个「小什么都不做」。

特别地，我们有一个对应图
\[
  \begin{tikzcd}
    & \Bun_G(\Sigma)\times_{\Bun_G(\Sigma\setminus x)}\Bun_G(\Sigma) \arrow[dl] \arrow[dr] & \\
    \Bun_G(\Sigma) & & \Bun_G(\Sigma).
  \end{tikzcd}
\]
我们用取各种层来线性化它，然后用推拉构造定义由某个局部算子诱导的映射。令 $\Hck_x := \Bun_G(\Sigma)\times_{\Bun_G(\Sigma\setminus x)}\Bun_G(\Sigma)$。也可以挖掉一个圆盘、粘进一个非 Hausdorff 的「加倍圆盘」来实现它：
\[
  \Hck_x \simeq \Bun_G(\Sigma)\times_{\Bun_G D_\varepsilon(x)}
  \bigl(\Bun_G(D_\varepsilon(x))\times_{\Bun_G(D_\varepsilon(x)\setminus x)}\Bun_G(D_\varepsilon(x))\bigr).
\]
于是我们得到：圆盘上有平凡化的主 $G$-丛空间，乘以圆盘上平凡 $G$-丛的修正。这把 $\Bun_G$ 以「adelic」方式实现为双商，正如我们在数域与函数域情形所做的。这个描述就是我们此前讨论过的 Hecke 对应：$\Hck_x$ 等同于
\[
  \Bun_G(\Sigma,x)\times\bigl(G(\cO_x)\backslash G(K_x)/G(\cO_x)\bigr).
\]
它是 $\Bun_G(\Sigma,x)$ 上重言 $G(\cO_x)$-丛所配的 $G(K_x)/G(\cO_x)$-丛的全空间。

物理上的要点是：$\cA_G(S^2)$，即 $\Bun_G(S^2)$ 上的层，通过线算子作用，而这里充满了有趣的算子。群论图像则是：起作用的是 $G(\cO_x)\backslash G(K_x)/G(\cO_x)$ 上的层，或意大利饺上的层。二者的关系是：如果你给意大利饺加足够多的馅，它看起来就像 $S^2$（不过我们要指出，无论加不加，意大利饺都很好吃）。

这个范畴，无论以哪种化身出现，都值得一个专门的名字。我们称它为球面 Hecke 范畴（spherical Hecke category），记作 $\mathrm{Sph}$。物理／TFT 的视角说：4 维理论中的线算子彼此交换，它们构成对称幺半范畴。群论视角对它们为何交换没有清楚的解释；而物理澄清了数学，即便在算术情形也是如此。
